\kbvsechead{6. ಈ ಪದಗಳ ಅರ್ಥ ವಿವರಣೆ ಕೊಡಿ :}

\vspace{2.5pt plus 3pt minus 0.5pt}\noindent\hspace*{1.2em}\kbvnum{(೧)}\hspace{0.55em}ಲಘು, (೨) ಗುರು\par

\vspace{2.5pt plus 3pt minus 0.5pt}\noindent\hspace*{1.2em}\kbvnum{7.}\hspace{0.55em}ಪ್ರಸ್ತಾರ ಹಾಕುವುದು ಎಂದರೇನು ?\par

\vspace{2.5pt plus 3pt minus 0.5pt}\noindent\hspace*{1.2em}\kbvnum{8.}\hspace{0.55em}ಈ ಕಂದ ಪದ್ಯಕ್ಕೆ ಪ್ರಸ್ತಾರ ಹಾಕಿ ಮಾತ್ರಾಗಣಗಳನ್ನು ವಿಂಗಡಿಸಿ.\par

\begin{kbvverse}
\noindent ಗಿರಿಜೆ ಶಶಿವದನೆ ಕೊಳನಂ । \\
\noindent ಭರದಿಂ ಪೊಗೆ ನೆಯ್ದಿಲೆಯ್ದಿದುವು ವಿಕಸನಮಂ । \\
\noindent ವಿರಹವನಾಂತವು ಕೋಳಂ । \\
\noindent ಸರಸಿರುಹಂ ಮುಗಿದುವರರೆ ಪಗಲೊಳ ಚಿತ್ರಂ ॥
\end{kbvverse}

\noindent\hfill{\fontsize{9.5pt}{12.5pt}\selectfont\color{kbvgrey}—ಹರೀಶ್ವರನ ಗಿರಿಜಾ ಕಲ್ಯಾಣ}\par\vspace{3pt plus 2pt}

\vspace{2.5pt plus 3pt minus 0.5pt}\noindent\hspace*{1.2em}\kbvnum{9.}\hspace{0.55em}ಮೇಲಿನ ಪದ್ಯವನ್ನು ಆಧಾರವಾಗಿಟ್ಟುಕೊಂಡು ಕಂದ ಪದ್ಯದ ಲಕ್ಷಣಗಳನ್ನು ತಿಳಿಸಿ.\par

\vspace{2.5pt plus 3pt minus 0.5pt}\noindent\hspace*{1.2em}\kbvnum{10.}\hspace{0.55em}ಷಟ್ಪದಿ ಎಂದರೇನು ? ಷಟ್ಪದಿಯಲ್ಲಿ ಎಷ್ಟು ಪ್ರಕಾರಗಳಿವೆ ? ಮತ್ತು ಅವು ಯಾವುವು ?\par

\vspace{2.5pt plus 3pt minus 0.5pt}\noindent\hspace*{1.2em}\kbvnum{11.}\hspace{0.55em}ನಿಮ್ಮ ಪಠ್ಯಪುಸ್ತಕದಲ್ಲಿರುವ ಷಟ್ಪದಿ ಜಾತಿಯ ಯಾವುದಾದ ರೊಂದು ಪದ್ಯಕ್ಕೆ ಪ್ರಸ್ತಾರ ಹಾಕಿ.\par

\vspace{2.5pt plus 3pt minus 0.5pt}\noindent\hspace*{1.2em}\kbvnum{12.}\hspace{0.55em}“ರಗಳೆ” ಎಂದರೇನು ? ಮುಖ್ಯವೆನಿಸಿದ ಮೂರು ಪ್ರಕಾರದ ರಗಳೆಗಳನ್ನು ಹೆಸರಿಸಿ.\par

\vspace{2.5pt plus 3pt minus 0.5pt}\noindent\hspace*{1.2em}\kbvnum{13.}\hspace{0.55em}ಲಲಿತ ರಗಳೆಯ ಲಕ್ಷಣಗಳನ್ನು ತಿಳಿಸಿ.\par

\vspace{2.5pt plus 3pt minus 0.5pt}\noindent\hspace*{1.2em}\kbvnum{14.}\hspace{0.55em}ವಾರ್ಧಕ ಷಟ್ಪದಿಯ ಲಕ್ಷಣಗಳನ್ನು ತಿಳಿಸಿರಿ.\par

\vspace{2.5pt plus 3pt minus 0.5pt}\noindent\hspace*{1.2em}\kbvnum{15.}\hspace{0.55em}ಯಾವ ಪ್ರಕಾರದ ಷಟ್ಪದಿ ಪದ್ಯಗಳಲ್ಲಿ ಮಧ್ಯೆ ಗುರುವುಳ್ಳ ಗಣ ಕಂಡುಬರುವುದಿಲ್ಲ ?\par

\vspace{2.5pt plus 3pt minus 0.5pt}\noindent\hspace*{1.2em}\kbvnum{16.}\hspace{0.55em}ಅಕ್ಷರ ಗಣ ಎಂದರೇನು ? ಅಕ್ಷರಗಣಕ್ಕೂ ಮಾತ್ರಾಗಣಕ್ಕೂ ಇರುವ ವ್ಯತ್ಯಾಸವೇನು ?\par

\vspace{2.5pt plus 3pt minus 0.5pt}\noindent\hspace*{1.2em}\kbvnum{17.}\hspace{0.55em}ವೃತ್ತಜಾತಿಯ ಛಂದಸ್ಸಿಗೆ ಸೇರಿದ ಆರು ವಿಧವಾದ ವೃತ್ತ ಗಳನ್ನು ಹೆಸರಿಸಿ.\par
